% SUPPLEMENT Table S2 since the 2026-09-23 six-figure wrap (formerly main Table 2); input by supplement.tex.
% Table 2 -- detection per held-out device and training arm at both operating points, with case-clustered intervals
% and PR-AUC. No merged cells: the target block heading sits in the first cell of its first row. Sources, one line per
% cell, in tab2_detection.sources.md; the same rows with BCa intervals, plus the arms cut from the page for space
% (the design-A seed pairs, renders-only on COBRA and DeepFluoro, Cios-only on COBRA, Elekta-only and Surgio-only on
% DeepFluoro), in tab2_detection.md.
%   results/a1.json                  the reported detector (lodo_carm) on C-ArmLSpine at 0.74 and 0.25 (k/n, percentile, BCa)
%   results/lcurve/lcurve.json       design-B stage-1 arms and the same-device controls at 0.74 and 0.25, PR-AUC (both devices)
%   results/contribution_matrix.json every other arm at 0.25 (recall as a fraction; k = recall x n, exact integers)
%   results/a3.json                  the plain-film block (render-trained vs film-trained, conf 0.05 and 0.25)
%   results/fig3_detection/scored/*  the 0.74 cells of the render-only and real-only held-out arms (same sweeps, both points)
% BOX CONVENTION (A28): every number here is in the whole-vertebra convention, and the sentence in the text that quotes a
% row must say "whole-vertebra". The vertebral-body-box row per block (the all-device detector det_g3_primary_s0, 15 G5
% cells) was DROPPED from paper 1 on 2026-09-24 (Greg): it can be scored only on the held-out test persons paper 2 keeps
% blinded; the rows are kept verbatim in paper2_levelcheck/MOVED_FROM_PAPER1.md (item 10).
% Labels: tab:detection is the table; tab:contribution is kept as an alias so the Results prose that still cites the
% retired contribution table resolves to this one (the single-source arms it quoted are rows here).
% Revision 2 (2026-09-24, REVISION_LOG2.md): plain-radiograph rows re-scored on the 17 NHANES II lumbar spines neither
% detector trained or was selected on (EVAL_LEAKS L2; closure MANUSCRIPT_CORRECTIONS.md 3a;
% results/eval_gate/leaks/render_vs_film_unseen_rerun.json); zero-floor percentile cells given unclustered Wilson
% intervals (REVIEW2_STATS N18; results/stats_review2/recompute.json#S10_floor_cells).
% Layout: landscape (rotated 23.2 cm minipage, see tab_corpus.tex); body 7 pt on 7.6 pt.
\begin{table}[p]
\centering
\begingroup
\singlespacing
\fontsize{7}{7.6}\selectfont
\newcommand{\fnm}[1]{\textsuperscript{\fontsize{7}{7}\selectfont #1}}
\setlength{\tabcolsep}{2pt}
\renewcommand{\arraystretch}{1.0}
\rotatebox[origin=c]{90}{\begin{minipage}{23.2cm}
\captionsetup{font=footnotesize,skip=4pt}
\centering
\caption{Detection per held-out device and training arm at the cross-device (0.74) and default (0.25) operating points. Whole-vertebra box convention throughout; the all-device detector trained on vertebral-body boxes is reported in the separate held-out study.}
\label{tab:detection}\label{tab:contribution}
% widths sum to 22.00 cm; 14 gaps of 2 pt (0.99 cm) -> 22.99 cm
\begin{tabular}{%
>{\raggedright\arraybackslash}p{2.20cm}%
>{\raggedright\arraybackslash}p{5.20cm}%
>{\raggedright\arraybackslash}p{3.15cm}%
>{\raggedright\arraybackslash}p{3.00cm}%
>{\raggedright\arraybackslash}p{3.15cm}%
>{\raggedright\arraybackslash}p{3.00cm}%
>{\raggedright\arraybackslash}p{2.30cm}}
\toprule
Target device\fnm{a} & Training arm & Vertebra recall at 0.74\fnm{b} & Sacrum found at 0.74 & Vertebra recall at 0.25 & Sacrum found at 0.25 & PR-AUC\fnm{c}, vertebra / sacrum \\
\midrule
C-ArmLSpine (Surgio-B51) & Reported detector: real COBRA2026 and DeepFluoro frames, no render & 80.5 (231/287) [74.1, 86.8] & 76.8 (43/56) [64.3, 89.3] & 90.6 (260/287) [84.3, 94.5]\fnm{g} & 83.9 (47/56) [73.2, 92.9] & -- \\
 & Design B: styled renders and real frames of the two other devices; seed 1 & 63.1 (181/287) [54.8, 71.0] & 66.1 (37/56) [53.6, 78.6] & 87.8 (252/287) [81.4, 93.2] & 85.7 (48/56) [76.8, 92.9] & 0.890 / 0.900 \\
 & Elekta XVI only, styled and real & -- & -- & 86.8 (249/287) [81.1, 91.8] & 78.6 (44/56) [67.9, 89.3] & -- \\
 & Cios Fusion only, styled and real & -- & -- & 12.9 (37/287) [8.3, 18.1] & AP 28/28, lateral 0/28\fnm{g} & -- \\
 & Elekta XVI only, real frames cut to 240 in 10 patients & -- & -- & 39.0 (112/287) [33.2, 44.7] & 37.5 (21/56) [28.6, 46.4] & -- \\
 & Renders only, no real frame & 0.3 (1/287) [0.1, 1.9]\fnm{g} & 5.4 (3/56) [1.8, 14.6]\fnm{g} & 31.7 (91/287) [27.1, 36.1] & 53.6 (30/56) [44.6, 62.5] & -- \\
 & Control\fnm{d}: C-ArmLSpine real training frames only & 40.4 (116/287) [37.8, 42.9] & 94.6 (53/56) [85.4, 98.2]\fnm{g} & 41.5 (119/287) [38.8, 44.0] & 98.2 (55/56) [90.6, 99.7]\fnm{g} & 0.367 / 0.979 \\
 & Control\fnm{d}: real training frames of all three devices & 40.1 (115/287) [37.7, 42.3] & 96.4 (54/56) [87.9, 99.0]\fnm{g} & 40.8 (117/287) [38.4, 43.0] & 98.2 (55/56) [90.6, 99.7]\fnm{g} & 0.385 / 0.982 \\
\addlinespace[1.5pt]
COBRA2026 (Elekta XVI) & Real DeepFluoro and C-ArmLSpine frames, no render & 6.2 (55/893) [2.9, 9.4] & 3.4 (21/609) [1.3, 5.9] & 18.5 (165/893) [13.7, 22.8] & 13.3 (81/609) [8.1, 18.8] & -- \\
 & Design B: styled renders and real frames of the two other devices & 10.1 (90/893) [6.7, 13.1] & 13.0 (79/609) [9.0, 17.1] & 32.7 (292/893) [28.4, 36.2] & 36.5 (222/609) [31.6, 40.9] & 0.266 / 0.391 \\
 & Surgio-B51 only, styled and real; two runs of one list & -- & -- & 7.5 (67/893) [3.5, 12.1]; 12.2 (109/893) [8.0, 16.0] & 3.8 (23/609) [1.8, 6.2]; 8.0 (49/609) [5.3, 10.9] & -- \\
 & Control\fnm{d}: real training frames of all three devices & 93.2 (832/893) [86.4, 97.7] & 90.6 (552/609) [83.8, 96.3] & 97.4 (870/893) [93.3, 99.8] & 97.9 (596/609) [95.7, 99.5] & 0.963 / 0.985 \\
\addlinespace[1.5pt]
DeepFluoro (Cios Fusion)\fnm{e} & Real COBRA2026 and C-ArmLSpine frames, no render; test specimen & 74.7 (62/83)\fnm{e} & 100.0 (90/90)\fnm{e} & 95.2 (79/83)\fnm{e} & 100.0 (90/90)\fnm{e} & -- \\
 & Design B: styled renders and real frames of the two other devices; six specimens & -- & -- & 78.6 (338/430) [68.8, 92.0] & 95.0 (321/338) [78.5, 99.8] & -- \\
 & The same arm; test specimen & 69.9 (58/83)\fnm{e} & 87.8 (79/90)\fnm{e} & 94.0 (78/83)\fnm{e} & 100.0 (90/90)\fnm{e} & 0.923 / 0.954 \\
 & Control\fnm{d}: real training frames of all three devices; test specimen & 73.5 (61/83)\fnm{e} & 100.0 (90/90)\fnm{e} & 95.2 (79/83)\fnm{e} & 100.0 (90/90)\fnm{e} & 0.685 / 0.968 \\
\addlinespace[1.5pt]
Plain radiographs\fnm{f} & Render-trained detector, no real image & at 0.05: 67.7 (109/161) & -- & 34.8 (56/161) & -- & -- \\
 & Film-trained detector & at 0.05: 99.4 (160/161) & -- & 98.8 (159/161) & -- & -- \\
\bottomrule
\end{tabular}
\par\vspace{1pt}
\begingroup\fontsize{7}{7.6}\selectfont\raggedright
Cells are \% ($k/n$) [95\,\% CI]. CI = confidence interval; PR-AUC = area under the precision--recall curve; IoU = intersection over union. Greedy one-to-one matching at IoU $\ge$ 0.5; 95\,\% percentile bootstraps of 2,000 resamples clustered by case (C-ArmLSpine), patient (COBRA2026) or specimen (DeepFluoro), BCa intervals in the Markdown twin; ``--'' = not scored at that point or not computed. Styled = rendered through the fitted terms; design B = renders under each device's own poses.
\fnm{a}Test sets: C-ArmLSpine 56 frames, 28 cases, 287 vertebrae, 56 sacra; COBRA2026 648 frames, 27 patients, 893 vertebrae, 609 frames with a sacrum; DeepFluoro six specimens, 366 frames, 430 vertebrae, 338 with a sacrum; plain radiographs 17 NHANES II lumbar films (17 spines, 161 bodies) whose patients neither detector trained or was selected on. Truth boxes: C-ArmLSpine rows against whole-vertebra outline boxes projected from the CBCT labels, the convention every number on disk uses (under the SPINEPS vertebral-body boxes of the pending flagship rows the reported detector reads 5.2\,\% (15/287) at 0.25, a convention mismatch); COBRA2026 and DeepFluoro rows against whole-vertebra boxes projected from the registered CT. The device is held out unless the row says control.
\fnm{b}0.74 = best pooled vertebra F1 of the reported detector's sweeps on COBRA2026, DeepFluoro and the Geneva phantom, reading no C-ArmLSpine frame, adopted after the 0.25 result was reported; 0.25 = the software default; the reported detector chose its checkpoint on the COBRA2026 and DeepFluoro test splits.
\fnm{c}Non-interpolated average precision over the confidence sweep, computed for the design-B arms and the controls only.
\fnm{d}Controls trained on the target device's own frames (test cases held out by case); on C-ArmLSpine the outline convention penalises these body-trained arms, which read 89.9\,\% (258/287) and 90.6\,\% (260/287) at 0.25 under body boxes.
\fnm{e}One specimen is one cluster, no interval; the six-specimen rows test arms that never saw the device, their six-cluster intervals descriptive.
\fnm{f}Scored at 0.05 and 0.25 only, on every annotated NHANES~II lumbar spine that neither detector trained or was selected on; paired difference at 0.05 +31.7 percentage points [21.0, 45.4], sign test 16 vs 0 spines, Holm-adjusted $P<.001$; 17 clusters, so the intervals are approximate.
\fnm{g}Wilson 95\,\% interval where the percentile one reached 0 or 100\,\% against the data (unclustered for the two 0 cells) or the file flags errors in few clusters (sacrum controls unadjusted, design effect about 1); Cios Fusion only: split by view, no pooled interval.
\par\endgroup
\end{minipage}}
\endgroup
\end{table}
